\ifdefined\MyWebBook\else
\documentclass[../textbook.tex]{subfiles}
\begin{document}
\fi
\bookchapter{检索增强生成}{ch:rag-systems}

RAG 系统将检索材料组织为生成上下文，并建立答案声明与来源的对应关系。证据选择需要处理重复、时效、例外条款与覆盖范围；知识更新和权限变化后，相关索引、引用与缓存也需同步维护。

\bookxref{ch:information-retrieval}讨论索引、词法与向量检索的基础。本章承接候选召回，进一步建立重排、证据选择、生成与知识服务的关系。

\section{检索增强的概率含义}

\subsection{检索作为潜变量}
\term{检索增强生成}{Retrieval-Augmented Generation}{RAG}利用外部材料改变生成条件。原始 RAG 将检索文档看作\term{潜变量}{Latent Variable}{}，结合检索分布与条件生成模型，而不只是把搜索结果放进提示词。\cite{lewis2020rag}

\begin{assumption}[概率模型与服务边界]
给定查询 $q$ 和固定知识快照 $\mathcal D$，检索对象 $z$ 为文档或片段，答案 $y=(y_1,\ldots,y_T)$ 为有限词元序列。概率模型中的文档分布已归一化。实际系统中，知识范围还受身份、权限和时间条件约束；所有生成器可见的证据必须先满足这些外部约束。检索分数、重排分数与答案正确概率不得未经校准混用。
\end{assumption}

本章主要符号见表\ref{tab:rag-systems-symbols}。
\begin{longtable}{@{}p{.32\linewidth}p{.64\linewidth}@{}}
\caption{检索增强生成的主要符号与形状}\label{tab:rag-systems-symbols}\\
\toprule 符号 & 含义与形状\\\midrule
\endfirsthead
\toprule 符号 & 含义与形状（续）\\\midrule
\endhead
\midrule
\multicolumn{2}{r}{续下页}\
\endfoot
\bottomrule
\endlastfoot
$q,y,z,\mathcal D$ & 查询、生成答案、潜在文档与知识快照。\\
$p_\eta(z\mid q),p_\theta(y\mid q,z)$ & 参数为 $\eta$ 的检索分布和参数为 $\theta$ 的条件生成分布。\\
$\mathcal C_K,E$ & 召回的 $K$ 个候选与最终有序证据包。\\
$f_\eta(q),g_\eta(d)\in\mathbb R^h$ & 查询与文档的 $h$ 维稠密表示。\\
$s_i,\gamma_i,l_i$ & 候选分数、给定答案后的文档后验责任、证据序列化词元成本。\\
$u_i\in\{0,1\},C_E$ & 证据选择变量与证据词元预算。\\
$A_{ji}$ & 证据 $i$ 是否覆盖问题要点 $j$ 的指示量。\\
$I,P,v$ & 已认证身份、当前权限策略、知识或组件版本。\\
\end{longtable}

对整条回答共享一个潜在文档，\term{序列级 RAG}{RAG-Sequence}{}写为
\begin{equation}
 p_{\mathrm{seq}}(y\mid q)
 =\sum_{z\in\mathcal D}p_\eta(z\mid q)
   \prod_{t=1}^{T}p_\theta(y_t\mid q,z,y_{<t}).
 \label{eq:rag-sequence}
\end{equation}
先在每个文档条件下计算整条答案概率，再对文档求和。对每个词元分别边缘化，\term{词元级 RAG}{RAG-Token}{}则为
\begin{equation}
 p_{\mathrm{token}}(y\mid q)
 =\prod_{t=1}^{T}\sum_{z\in\mathcal D}p_\eta(z\mid q)
  p_\theta(y_t\mid q,z,y_{<t}).
 \label{eq:rag-token}
\end{equation}
求和与乘积位置不同，意味着模型假设不同：后者允许不同生成位置由不同文档项支持。\keyconcept{序列级 RAG 假定整条回答由单一潜在文档完全支持，而词元级 RAG 允许不同生成位置由不同文档动态支持；两者的求和与连乘次序反映了本质不同的概率建模假设。}这里的词元级潜变量边缘化是概率模型内部的结构，与服务端每生成一个词元是否重新调用搜索引擎没有关系。

大语料无法直接遍历全部文档，通常只计算候选集 $\mathcal C_K$。如果保留原完整分布的概率，截断求和是一种近似，其总质量可能小于1；如果在候选集重新归一化，定义的是条件于候选集合的分布。记原候选质量为 $Z_K=\sum_{z\in\mathcal C_K}p_\eta(z\mid q)$，则局部归一化使用 $\widetilde p_\eta=\frac{p_\eta}{Z_K}$。两种口径在训练与解释时必须明确区分。

\begin{example}


设两个候选文档权重为 $(0.6,0.4)$，目标答案由两个词元组成。第一文档对两个词元的条件概率分别为0.9和0.2，第二文档分别为0.1和0.8；第二个概率均已条件于相同的目标前缀。序列级概率为
\begin{equation*}
 p_{\mathrm{seq}}=0.6(0.9\cdot0.2)+0.4(0.1\cdot0.8)=0.14.
\end{equation*}
词元级概率则为
\begin{equation*}
 p_{\mathrm{token}}=(0.6\cdot0.9+0.4\cdot0.1)
 (0.6\cdot0.2+0.4\cdot0.8)=0.2552.
\end{equation*}
词元级模型在不同位置混合文档贡献，因此两种概率形式得到不同结果。数值描述给定序列的生成概率，事实正确性通过任务标准评价。

序列级模型还可推导检索训练信号。令 $a_i=p_\theta(y\mid q,z_i)$，$p_i=\softmax(s)_i$，$Z=\sum_i p_i a_i$。给定答案后的文档后验为
\begin{equation*}
 \gamma_i=\frac{p_i a_i}{Z}.
\end{equation*}
利用 softmax 导数 $\frac{\partial p_i}{\partial s_j}=p_i(\mathbf1[i=j]-p_j)$，
\begin{align*}
 \frac{\partial Z}{\partial s_j}&=p_ja_j-p_jZ,\\
 \frac{\partial(-\log Z)}{\partial s_j}&=p_j-\gamma_j.
\end{align*}
例题中 $\gamma_1=\frac{0.108}{0.14}=\frac{27}{35}\approx0.7714$，$\gamma_2\approx0.2286$。第一文档的梯度约为 $-0.1714$，梯度下降会提高其检索分数。该梯度提高有助于解释训练答案的候选分数，直接更新范围限于候选集合。文档真实性与引用充分性由相应标注和检查确定。
\end{example}

\subsection{检索上下文组装}
常见服务先选择并序列化证据包 $E=\mathcal A(q,\mathcal C_K)$，再调用
\begin{equation}
 p_\theta(y\mid q,\operatorname{serialize}(E)).
\end{equation}
这允许生成器同时阅读多个片段并综合条件，式\eqref{eq:rag-sequence}中的文档概率边缘化在这里由注意力隐式承担。重排分数只决定排序或选择，除非目标另行定义，并不会自动成为生成概率中的混合权重。

证据拼接具有明确的工程可控性：可以带上来源、时间、权限与引用标识，也容易替换检索器。生成器的注意力分配描述内部聚合过程，声明与来源的支持关系通过内容核验建立。\footnote{当前工程语境常把多种“检索后生成”架构统称为 RAG。本章保留这一广义用法，同时用 RAG-Sequence 与 RAG-Token 指称原论文中的具体概率形式，避免同名导致目标混淆。}

\section{召回、重排与训练关系}

\subsection{多路检索的目标}
\term{稀疏检索}{Sparse Retrieval}{}通过词项维度匹配材料，BM25 使用词频饱和、逆文档频率和长度校正；\term{稠密检索}{Dense Retrieval}{}通过学习表示的相似度选择材料。基础公式见前章。本章关注它们在管线中的不同职责：词法路径保留编号、专名和少见术语的精确线索，语义路径连接不同表达；完整限定的覆盖率由目标查询集上的召回评估确定。

稠密双编码器常使用 $s(q,d)=f_\eta(q)^{\mathsf{T}} g_\eta(d)$。对正例 $d^+$ 与候选负例集 $\mathcal N$，一种训练目标为
\begin{equation}
 \mathcal L_{\mathrm{ret}}=-\log
 \frac{\conste{}^{\frac{s(q,d^+)}{\tau_r}}}
 {\conste{}^{\frac{s(q,d^+)}{\tau_r}}+\sum_{d^-\in\mathcal N}\conste{}^{\frac{s(q,d^-)}{\tau_r}}},
\end{equation}
其中 $\tau_r>0$ 是排序温度。它使正例相对负例更靠前；负例采样和潜在假负例会改变学习信号。DPR 是采用双编码器进行段落检索的重要实例。\cite{karpukhin2020dpr} 相对分数用于候选排序。

\term{混合检索}{Hybrid Retrieval}{}将多路候选取并集，再做融合。若使用分数线性组合 $s_h=\alpha\widetilde s_{\mathrm{lex}}+(1-\alpha)\widetilde s_{\mathrm{dense}}$，应先定义两路尺度校准；原始 BM25 与余弦分数直接相加缺少稳定的数值解释。也可按名次融合，例如 $s_h(d)=\sum_j\frac{1}{c+\operatorname{rank}_j(d)}$，其中 $c>0$ 平滑名次影响，未召回项不贡献分数。此类融合得到的是排序量，与经过概率归一化的证据置信度属于不同量纲。

并集召回提高候选覆盖的可能性，也会增加重复、噪声和后续重排成本。只有被所有查询路径统一授权的候选才能进入并集，任何一条备用路径都受同样的身份过滤约束。

\subsection{联合编码重排}
\term{重排}{Re-ranking}{}在较小候选集上计算更精细的查询与文档关系。\term{交叉编码器}{Cross-encoder}{}共同编码查询和候选内容，可以直接比较限定条件、否定关系和实体对应，通常比独立向量内积具有更强的交互能力，但每个查询文档对都需执行模型。使用 BERT 进行段落重排的工作展示了这一模式。\cite{nogueira2019reranking}

具体模型如 \texttt{BAAI/bge-reranker-v2-m3}：将原始问题与每个召回片段组成输入对，模型输出用于候选排序的相关性分数\footnote{BAAI，\texttt{bge-reranker-v2-m3} 模型卡，版本 \texttt{953dc6f}，\url{https://huggingface.co/BAAI/bge-reranker-v2-m3/blob/953dc6f6f85a1b2dbfca4c34a2796e7dde08d41e/README.md}，访问日期：2026-09-23。}。可用 \texttt{bge-m3} 生成向量进行首轮语义召回，再由重排模型比较查询和候选正文；两者也可以与 BM25 等其他召回器组合。重排器输出只在相同模型、输入协议和候选集合的排序任务下解释，既不自动给出声明所需的证据支持，也不能直接充当拒答概率。每对输入都要在线计算，候选数与输入截断长度应同召回覆盖和时延一起测量。

若候选得分为 $a_i=r_\phi(q,d_i)$，成对目标可写为
\begin{equation}
 \mathcal L_{\mathrm{pair}}=-\log\sigma(a_+-a_-).
\end{equation}
其对正例分数的导数为 $-\sigma(a_--a_+)$，对负例为相反数，促使正例相对提升。若有多个相关候选，列表目标可用标签分布 $t_i$ 对 softmax 分数做交叉熵：
\begin{equation}
 \mathcal L_{\mathrm{list}}=-\sum_i t_i\log\frac{\conste{}^{a_i}}{\sum_j \conste{}^{a_j}}.
\end{equation}
成对和列表目标都必须明确相关性含义。与问题主题相似、包含某个答案词、足以支持结论，分别对应不同标签；训练标注主题相关性而部署要求证据充分性，会形成目标错位。

\emph{重排器仅能调整候选集内部的相对次序}，对第一阶段多路召回中\emph{彻底漏检的文档无能为力}；因此召回阶段的覆盖率决定了 RAG 系统的性能上限。候选文本截断还可能丢失例外条款，使重排器看见的内容与生成器最终读取的材料不同。因此重排输入范围和证据正文范围应可追踪，浮点分数之外还需要记录实际文本范围。

\subsection{检索模型训练}
嵌入模型优化召回，重排器优化候选顺序，生成器优化在证据条件下回答。三者可以分别训练，也可让后者提供监督，但训练依赖必须明确。前例中似然对检索分数的梯度 $p_j-\gamma_j$ 提供一种候选监督信号；离散近邻搜索的候选成员变化在普通反向传播中不可微，因此“联合训练”通常还是回归到固定候选、近似梯度或交替刷新索引。

文档编码器更新后，旧向量索引与查询向量不再同属一致表示空间，必须重建或采用明确的过渡策略。只更新查询侧、冻结文档编码器可以降低刷新成本，代价是可调整范围收窄。重排或生成模型升级也可能改变证据选择和引用行为，即使文档索引未变仍需独立比较。

生成器训练需要包含充分证据、缺失证据和冲突证据，全部使用完美材料的训练集会鼓励始终作答，部署时的拒答就只能依赖脆弱的提示。用模型自身答案生成检索监督还可能强化既有错误，教师或生成分数应视为弱监督并保留独立证据标注。

\begin{tip}[拒答与鲁棒性训练中的负样本配比]
RAG 生成器微调语料必须包含“证据充分可回答”、“证据不足需弃权”以及“证据相互冲突”三种真实场景。若训练集全部由完美支持的文档构成，模型将过度拟合“强行从上下文中寻找答案”的先验，丧失必要的拒答能力与事实审慎性。
\end{tip}

\section{证据选择}

\subsection{有效长度约束}
\term{证据组装}{Evidence Assembly}{}把候选变成可被生成器读取的有序材料。总上下文容量 $C_{\mathrm{ctx}}$ 应满足
\begin{equation}
 C_E\leq C_{\mathrm{ctx}}-C_{\mathrm{sys}}-C_q-C_{\mathrm{hist}}-C_{\mathrm{out}}-C_{\mathrm{reserve}}.
\end{equation}
各项分别是系统模板、问题、对话历史、输出与保留余量。证据成本 $l_i$ 必须按实际分词器和序列化格式计算，包含标题、来源与引用标识，而不只是正文字符数。预算为负时应调整上下文策略，静默裁掉必要约束会使约束条件消失于生成器视野。

选择变量 $u_i\in\{0,1\}$ 满足 $\sum_i u_il_i\leq C_E$。最大化 $\sum_i u_i s_i$ 只优化带预算的相关性；多要点覆盖与去重需要加入相应目标。若问题被分为要点 $j$，覆盖指示为 $A_{ji}$，可定义
\begin{equation}
 F(S)=\sum_j w_j\min\left(1,\sum_{i\in S}A_{ji}\right)
 -\lambda\sum_{i<k,\ i,k\in S}r_{ik},
\end{equation}
其中 $r_{ik}\geq0$ 衡量重复，$w_j$ 表示要点权重。覆盖部分具有重复收益递减的性质；加入任意重复惩罚后，整体随集合增大的单调性消失，简单贪心的最优性保证也需重新判断。

来源权限、有效日期和必要限定通常应作为约束，用一个小的罚分抵消会允许它们被相关性压过。例如证据只在某版本产品下适用，把版本不符折价成扣分就等于承认它可以入选。多个来源互相转述属于同一条证据链，来源多样性应追踪原始出处。

\begin{example}


设总上下文为1024词元，系统模板180、问题80、历史160、输出240、余量44，可用证据预算为320。查询需要三个要点：A为适用条件，B为例外情况，C为生效日期。候选如下，长度已包含引用元数据。
\begin{table}[htbp]
\centering
\caption{多要点证据选择特征与覆盖关系}
\label{tab:rag-example-candidates}
\begin{tabular}{crcl}
\toprule
片段 & 长度 & 重排分数 & 支持要点\\
\midrule
$d_1$ &220&0.95&A\\
$d_2$ &180&0.85&A、B\\
$d_3$ &120&0.75&C\\
$d_4$ &100&0.70&B\\
\bottomrule
\end{tabular}
\end{table}
顺序装入最高分 $d_1$ 后只剩100词元，可以装入 $d_4$，用满320却缺少C。选择 $d_2,d_3$ 只用300词元，覆盖全部三个要点。即使采用分数总和，前一方案为1.65、后一方案为1.60，仍会偏向不完整的证据包。任务覆盖目标可优先选择包含全部必要要点的证据包。

如果 $d_2$ 的例外条款依赖前一段定义，应把必要上下文的额外成本纳入 $l_2$ 并重新选择；截断后仍声称覆盖B会高估证据完整度。若C的来源已过期，则两个方案都可能不足；预算优化只能在有效证据中选择，缺失事实不会因优化而出现。
\end{example}

\subsection{压缩、排序与引用锚点}
证据过长时，可以提取相关句段或生成摘要。\term{上下文压缩}{Context Compression}{}应保留数值、否定、范围与例外；生成摘要可能加入原文不存在的连接推断，因此\emph{摘要仍需映射回原始片段}。引用不应只指向一个不可复原的摘要缓存。

排序影响模型读取，但不是数学上的事实优先级。可按子问题组织证据，把一般规则与例外相邻放置，并标明来源日期和适用条件。重复块可合并展示，同时保留全部来源映射；多块共享同一引用标识会使声明定位模糊。证据 ID 应稳定指向文档版本及局部位置，随排序变化的“第3条”在下一次重排后就是另一个对象。


\section{检索策略}
\optional
\subsection{查询改写约束}
\term{查询改写}{Query Rewriting}{}把用户措辞转换为更适合检索的表达，例如补齐多轮对话中的指代、展开缩略语、产生同义表达。改写器可由规则或语言模型实现，也可根据后续阅读器的反馈训练\cite{ma2023queryrewrite}。面向搜索引擎的查询扩展与用户原意等价性存在张力：例如将“去年生效的规则”改写为“最新规则”，虽能提高某些高频文本的召回率，却抹除了时间边界这一关键约束条件。

查询状态因此至少分为三层：原始问题、可变的检索表达，以及只由控制平面维护的约束。用户身份、租户、授权范围、时间口径、已明确的实体标识和输出语言应由控制平面保存；每次改写只生成候选表达，原始请求保持不变。若对话指代仍有歧义，系统可以并行检索明确列出的候选实体，或要求澄清。把不确定指代强行绑定到最高频实体，后续证据会呈现表面一致而实际答非所问的结果。

设原始问题为 $q$，固定约束为 $b$，第 $j$ 个检索表达为 $q^{(j)}=g_j(q,b)$。多路检索取得
\begin{equation}
 \mathcal C=\bigcup_j\operatorname{Retrieve}(q^{(j)};\mathcal D_{I,P},b),
\end{equation}
其中 $\mathcal D_{I,P}$ 是当前身份和策略允许读取的语料。合并后按文档版本及片段位置去重，再使用原始问题和约束重排。这样可以避免改写器生成的一条宽泛查询主导最终答案，也使“某个候选来自哪条改写”成为可审计信息。

\subsection{分解与结构化知识}
\term{查询分解}{Query Decomposition}{}把复合问题拆为若干子问题。独立子问题可以并行；存在实体或时间依赖的子问题则应构成有向无环图。例如“某项目负责人管理的服务在上月发生了几次故障”，应先确定指定时点的负责人和服务范围，再查询相同时间窗口内的故障记录。按句号拆分只得到表面片段，依赖图还需编码实体与时间关系；现任负责人和上月负责人属于两个时点的状态。

结构化数据适合精确筛选、连接与聚合，文本适合解释规则与例外。系统可以把问题映射为受约束的查询计划，但字段、连接路径、读写权限、扫描范围和结果上限由执行层校验。模型提出的查询方案还要经过执行层授权。聚合结果应附带查询口径、数据快照、过滤条件和可追溯的来源记录范围；“总数为0”“没有匹配数据”和“权限不足，无法查询”必须使用不同状态，否则生成器会把三者都读成零。

图结构知识同样需要语义约束。关系能否沿路径传递取决于各边的语义及适用的推理规则；路径检索所得事实必须保留边类型、方向、有效时间与来源。把图路径和表格结果转换为统一证据对象，有助于生成器跨来源组织回答；统一格式的同时要保留原来的精确查询条件。

\subsection{线性、条件、分支与迭代流程}
线性流程按改写、检索、重排、组装、生成依次运行，适用于证据需求已知的问题。条件流程在确定性信息充分时选择不同路径，例如精确文档 ID 走版本读取，统计问题走结构化查询，解释性问题走文本检索。路由器的错误率本身需要评估，端到端失败在检索器之外还有路由器这一层来源。

分支流程同时调用不同检索通道，随后融合。设分支耗时为 $t_1,\ldots,t_J$，在资源充足且没有排队竞争的理想条件下，等待全部分支的耗时约为 $\max_jt_j+t_{\mathrm{merge}}$，而调用成本仍接近各分支成本之和。超时返回部分结果时，必须把缺失通道记录为未完成；记为“该通道没有证据”会把超时与空结果混为一谈。

\term{迭代检索}{Iterative Retrieval}{}根据当前证据缺口提出新查询。主动检索研究还探讨在生成过程中触发额外检索\cite{jiang2023activerag}。这里需要区分词元低概率和事实缺失：罕见专有名词可能概率低但证据充分，常见却错误的说法也可能概率高，两者在概率上不区分。触发信号决定何时检索，事实正确性由来源支持与后续核验判断。

令 $E_t$ 为第 $t$ 轮有效证据，$U_t$ 为尚未解决的事实需求集合，则一种明确的状态转移为
\begin{align}
 q_{t+1}&=\operatorname{Rewrite}(q,b,E_t,U_t),\\
 E_{t+1}&=\operatorname{Select}\bigl(E_t\cup\operatorname{Retrieve}(q_{t+1}),C_E\bigr),\\
 U_{t+1}&=\operatorname{Assess}(q,b,E_{t+1}).
\end{align}
式中的检索与选择均隐含授权和版本检查。需求评估可以由规则与模型共同完成，但必须保留不确定状态。停止条件至少包括证据已足够、连续一轮没有新增有效证据、迭代次数达到上限、词元预算耗尽和截止时间到达。终止后未解决的问题应进入有限回答或拒答路径；“轮数已经用完”是预算状态，与结论是否成立无关。

\begin{figure}[htbp]
\centering
\begin{tikzpicture}[node distance=7mm,box/.style={draw,rounded corners,align=center,text width=28mm,minimum height=10mm,font=\small},>=stealth]
\node[box] (q) {原始问题\\固定约束};
\node[box,right=of q] (r) {改写与分解\\检索计划};
\node[box,right=of r] (s) {多路检索\\授权过滤};
\node[box,below=of s] (e) {重排与证据选择\\预算检查};
\node[box,left=of e] (g) {生成与声明检查\\引用、冲突、拒答};
\node[box,left=of g] (o) {回答与来源\\运行清单};
\draw[->] (q)--(r); \draw[->] (r)--(s); \draw[->] (s)--(e);
\draw[->] (e)--(g); \draw[->] (g)--(o);
\draw[->] (e.east)--++(8mm,0)|-node[pos=.65,right,font=\small]{证据缺口，受限重试}(s.east);
\end{tikzpicture}
\caption[受约束的检索生成流程]{受约束的检索生成流程。迭代补充的是证据；身份、时间范围和资源上限由外部控制层保持。}\label{fig:rag-query-flow}
\end{figure}

\begin{algorithm}[有界证据检索与回答]
\begin{enumerate}
 \item 接收问题 $q$、身份 $I$、固定约束 $b$ 与截止时间；取得语料快照和当前授权策略。初始化证据 $E=\varnothing$，未解决需求 $U$。
 \item 在轮数、调用量、词元与时间均未超限时，依据 $q,b,U$ 生成检索计划；校验计划只包含允许的数据源和只读操作。
 \item 并行执行可独立的分支，区分成功为空、失败和超时。合并候选，检查当前权限、来源版本与有效时间，记录各分支来源。
 \item 用原始问题重排候选；在实际序列化词元预算内选择证据，保留定义、例外与引用锚点。更新 $U$。
 \item 若证据充分则结束检索；若没有新增有效证据或预算耗尽则保留未解决状态并结束；否则继续针对缺口检索。
 \item 根据证据生成有限范围的回答，把声明映射至来源。检查引用解析、语义支持和冲突；失败声明删除、限定或转为拒答。
 \item 输出前再次检查访问权限，返回答案状态、来源版本、未解决事项和运行清单。权限撤销使相应证据失效时，不返回依赖它的答案。
\end{enumerate}
\end{algorithm}

\section{基于证据的回答}

\subsection{声明证据映射}
\term{证据忠实度}{Evidence Faithfulness}{}衡量回答声明是否受到给定证据支持。答案正确性还取决于证据是否准确、有效，以及声明是否符合实际情况。例如，模型可能凭参数记忆答对，却引用了不相关段落；也可能准确复述一份已经过期的文件。因此，引用质量与答案正确性需要分别评估\cite{gao2023citations}。

设答案包含可验证声明 $c_1,\ldots,c_M$，来源集合为 $E$。引用关系 $J\subseteq\{1,\ldots,M\}\times E$ 不只是格式标记，还应附带支持、矛盾或证据不足的判定。检查分为三层：引用标识是否属于本次证据包；标识是否仍能解析到指定版本及片段；该片段或明确组合的片段是否支持声明。三层检查分别确定标识合法、来源可解析与内容支持。

一句话可能同时包含一般规则、时间条件与例外，应按可独立核实的主张划分；按标点机械切分会把同一条主张拆散或把多条主张并在一起。跨文档推导需标明各来源承担的前提；计算结论需保留输入数值、单位和计算规则。来源提到“多数情况适用”只支持多数情况；在规范性文本中，“建议”与“必须”具有截然不同的效力约束，严禁在生成中擅自提升约束强度。自动蕴含判别可作为检查器，但其输出依然需要在领域样本上校准，其误判率应通过领域标注估计。

\subsection{证据时效与冲突消解}
\term{数据新鲜度}{Data Freshness}{}描述服务使用的数据与所需时间状态之间的差距。抓取时间、发布日期和生效时间是不同字段。今天重新抓取的旧文件不会因此成为新规则；下月才生效的文件用于回答今天的状态也不合适。

处理冲突时，先确认实体、辖域、版本、有效时间和适用对象是否一致；若适用条件不同，应分别回答。只有条件相同而结论不一致时，才属于待解决的实质冲突。来源权威性也应由领域规则确定，文风和搜索名次只是间接线索。无法确定替代关系时，应并列说明冲突及来源时间，并收缩结论范围。

设某系统运行制度包含以下条款：版本 v1 规定普通批处理任务运行时长上限为 $\qty{30}{\minute}$，版本 v2 自特定生效日起调整为 $\qty{45}{\minute}$；同时另有特别条款声明“例行维护窗口内任务上限为 $\qty{10}{\minute}$”。若查询普通工作日的当前配额限制，须核验版本 v2 在查询时点的生效状态；若查询维护窗口的配额，须匹配 $\qty{10}{\minute}$ 的特例条款；若回溯历史日期，则须锚定该时点对应的版本 v1。将不同生效区间与特例条款的数值机械平均为 $\qty{28.3}{\minute}$ 毫无实际语义。若维护窗口的适用条件未在上下文中出现，系统须显式指出该信息缺口，直接输出检索分数最高的 $\qty{45}{\minute}$ 将严重误导调用方。

\subsection{选择性回答的决策条件}
\term{拒答}{Abstention}{}是证据不足时的一种合法输出状态，包括完全拒答和只回答已证实部分。访问拒绝、服务异常和问题超出支持范围属于另外三类状态，需要分别记录，否则质量指标会掩盖系统故障。

可以通过一个简化决策模型说明阈值的来源。设 $p$ 是在当前证据条件下某一拟答结论正确且适用的校准概率，答错成本为 $C_w$，拒答成本为 $C_r$，答对成本为0，且 $0<C_r<C_w$。回答与拒答的期望成本分别为
\begin{equation}
 \begin{aligned}
  R_{\mathrm{answer}} &= (1-p)C_w, \\
  R_{\mathrm{abstain}} &= C_r.
 \end{aligned}
\end{equation}
选择回答的条件为
\begin{equation}
 (1-p)C_w\le C_r
 \quad\Longleftrightarrow\quad
 p\ge1-\frac{C_r}{C_w}.
\end{equation}
例如当 $C_w=20, C_r=1$ 时，决策阈值为 $p \ge 1 - \frac{1}{20} = 0.95$。该临界值随错误惩罚与拒答成本的相对权重动态转移。真实系统还可能区分不同错误类型以及部分回答的成本。尤其不能把 BM25 分数、向量内积或语言模型最大词元概率直接当作 $p$ 的取值；这些量不具备所需的概率语义。

阈值应在与目标分布相符的开发集上选择，报告回答覆盖率和已回答问题的错误率，并单独分析知识库不可回答的问题。只奖励拒答准确率会诱导系统拒绝所有问题，只奖励已回答准确率又会诱导通过少答取得表面提升。可用性与错误风险应同时呈现。

\section{知识服务接口}

\subsection{离线知识管线}
完整服务首先需要可追溯的知识摄取流程：从授权来源取得原始文件及版本，解析为段落、表格和结构块，保存来源位置，按语义关系切分，计算检索表示，构建索引，最后发布可查询快照。解析与切分会改变信息可见性，必须保留原始文件到片段的映射。表头、单位、脚注与跨页续表若被丢弃，后续向量检索再准确也找不回丢失的语义。

\term{数据契约}{Data Contract}{}规定各阶段对字段、版本和失败状态的共同解释。文档对象至少包含稳定文档 ID、源版本、内容哈希、来源位置、有效时间、观测时间、授权标签和删除状态；片段额外包含父文档版本、局部位置、解析器与切分规则版本。嵌入记录还要包含模型和预处理版本。向量维度相同绝非表示空间兼容：模型权重的重训或升级会彻底改变流形几何，禁止因维度一致而混用不同版本嵌入模型生成的向量索引。

\term{快照清单}{Snapshot Manifest}{}记录一次发布使用的文档集合、索引分片、解析和嵌入版本，以及构建完成状态。新快照在暂存区完成后再切换查询指针，使一次查询使用可解释的版本组合。若数据量要求增量发布，也应显式定义可见水位及各分片版本；“最终会同步”描述的是收敛目标而非当前一致性。

\begin{figure}[htbp]
\centering
\begin{tikzpicture}[node distance=6mm,box/.style={draw,align=center,text width=27mm,minimum height=11mm,font=\small},>=stealth]
\node[box] (a) {授权原始来源\\版本、时间、权限};
\node[box,right=of a] (b) {解析与切分\\原始位置映射};
\node[box,right=of b] (c) {稀疏与稠密索引\\模型版本};
\node[box,below=of c] (d) {发布快照清单\\切换可见版本};
\node[box,left=of d] (e) {在线证据读取\\当前授权检查};
\node[box,left=of e] (f) {回答及缓存\\来源依赖};
\draw[->] (a)--(b);\draw[->] (b)--(c);\draw[->] (c)--(d);\draw[->] (d)--(e);\draw[->] (e)--(f);
\draw[->,dashed] (a.south)--(f.north) node[midway,left,font=\small]{撤销、删除};
\end{tikzpicture}
\caption[知识发布与在线读取边界关系]{知识发布与在线读取边界关系。内容快照保证可复现性，权限撤销仍须及时作用于在线访问。}\label{fig:rag-ingestion}
\end{figure}

\subsection{知识索引维护}
知识摄取常由\term{变更数据捕获}{Change Data Capture}{CDC}或周期扫描驱动。系统应假定事件可能重复、乱序或暂时丢失，并使用源版本与操作类型构造幂等键。重复处理同一版本不应生成新的逻辑片段，较旧更新覆盖较新的删除标记同样会破坏状态。没有全局递增版本的来源，可使用来源提供的条件读取与内容哈希，并明确它能判断的一致性关系仅限于其中一部分。

解析失败应保留原始对象和错误原因，便于使用修复后的解析器重放。嵌入失败可以重试该版本的片段；索引写入部分成功时，发布状态仍应保持未完成。重试次数耗尽的对象进入待修复队列，并纳入覆盖率和新鲜度报告。将失败文件悄悄丢弃，会使回答系统把摄取故障误认为知识不存在。

删除必须传播到片段、索引、摘要、引用解析服务和答案缓存。逻辑删除标记可先阻止在线读取，后台再回收物理数据；权限撤销则应优先改变访问判定，不等待所有索引副本重建。历史审计需要保留什么数据、保留多久，由独立的数据治理规则决定；“为了可复现”不是无限保留已失效内容的理由。

\subsection{在线对象与失败状态}
表\ref{tab:rag-contract}列出模块之间需要交换的核心对象。分数、文本与来源必须一起传递；只返回字符串列表会丢弃授权、版本和重试语义。

\begin{table}[htbp]
\centering
\caption[知识服务逻辑数据契约]{知识服务逻辑数据契约。具体编码形式可变，字段语义应保持稳定。}\label{tab:rag-contract}
\begin{tabularx}{\linewidth}{@{}lX@{}}
\toprule
对象 & 关键字段及含义 \\
\midrule
查询请求 & 原始问题、身份、租户、时间口径、截止时间、允许的数据源与输出约束 \\
候选证据 & 文档及片段 ID、源版本、通道及原始分数、位置、授权标签、有效时间 \\
证据包 & 已选择的原文、引用锚点、序列化词元数、覆盖需求、未解决冲突、快照 ID \\
回答对象 & 回答状态、文本、声明到来源映射、未解决需求、访问和引用检查结果 \\
运行清单 & 各模型与提示版本、检索参数、快照、分支状态、耗时、预算消耗和错误码 \\
\bottomrule
\end{tabularx}
\end{table}

在线状态应区分无相关候选、证据不足、来源冲突、无权限、上游不可用、预算耗尽和正常完成。某个非必要分支失败时可以在限定范围内继续；必要证据源不可用时应返回降级状态，把不完整证据包装成完整答案会掩盖缺口。生成阶段重试可复用同一证据快照，但需要再次读取当前权限。取消请求后，也应取消尚未执行的分支，避免孤立的检索与模型调用继续消耗资源。

流式输出会使检查更困难：一旦错误声明已经发给用户，后台撤回与从未发布并不等价。对需要强引用约束的答案，可先生成并检查声明片段再发送，或明确区分尚未核实的过程状态与最终答案。选择粒度会影响首字延迟，应作为服务契约的一部分。

\section{知识访问控制}

\subsection{权限过滤}
设授权判定为 $\operatorname{Allow}(I,P,d)$，可访问语料定义为
\begin{equation}
 \mathcal D_{I,P}=\{d\in\mathcal D:\operatorname{Allow}(I,P,d)=1\}.
\end{equation}
检索、重排、摘要和生成原则上都只能处理这个集合中的内容。只在最终答案中隐藏文档链接并不足够，未授权内容可能已经进入模型上下文、日志或缓存。向外部服务发送查询和证据，还需要满足相应数据处理边界。

索引能够执行权限预过滤时，应把身份约束纳入候选生成；如果某种近似索引只能在召回后过滤，必须评估过滤导致的候选不足，并确保未授权文本不进入后续模块。授权过滤前后的候选数量属于可能透露隐藏资源存在性的信息，其对外暴露范围应符合系统既定的资源存在性策略。

\term{提示注入}{Prompt Injection}{}在 RAG 中常来自检索文档内伪装成指令的文字。结构化分隔和明确提示有助于模型识别证据边界，但它们不是授权机制。租户、权限、数据源白名单和工具能力都属于模型外部的控制，检索内容对它们没有修改作用。审计日志本身也应受访问与保留策略约束，以免成为敏感证据的旁路副本。

\subsection{缓存一致性}
缓存减少重复计算，但回答缓存保存的是依赖特定知识状态和访问条件的结果。一个逻辑缓存键可以写为
\begin{equation}
 K_{\mathrm{cache}}=H(q_{\mathrm{norm}},I_{\mathrm{scope}},P_{\mathrm{epoch}},v_{\mathcal D},
 v_{\mathrm{models}},v_{\mathrm{prompt}},b,v_{\mathrm{config}}),
\end{equation}
其中 $H$ 是键构造函数，$I_{\mathrm{scope}}$ 是授权等价范围的标识，$P_{\mathrm{epoch}}$ 为策略版本。问题归一化不能删除否定、日期或实体限定。两个用户问题相同并不足以共享答案，还需要授权等价。

缓存还应记录来源 ID 和版本依赖。更新或撤销某个来源后，可使依赖该来源的答案失效；但新文档可能改变从未引用它的旧答案，尤其是原先“未找到证据”的\term{负缓存}{Negative Cache}{}。仅按已有引用做反向失效因此并不充分，还需要语料代际、主题分区版本或有明确定义的新鲜度窗口。

\term{存活时间}{Time to Live}{TTL}限制缓存陈旧的最长时间，权限检查仍需独立进行。权限撤销后，即使缓存尚未到期，也应阻止读取依赖被撤销证据的内容。另一方面，把每次摄取都绑定到全局缓存清空虽然简单，代价是集中失效和负载突增。可使用分区版本及带抖动的过期时间改善负载，但必须保留正确的依赖语义。

\begin{caution}[证据版本与当前授权]
内容快照说明“答案依据了什么版本”，当前授权说明“此刻是否允许读取这些内容”。权限撤销立即约束历史快照的读取；即使权限相同，缓存中的知识仍需满足问题要求的时间范围。
\end{caution}

\section{RAG 评估}

\subsection{检索生成质量链}
评估应使用保留原始问题、参考事实、可接受来源与时间范围的数据集，并记录知识库是否实际包含所需证据。对于可回答问题，检索层评估相关片段召回与排序；证据组装层评估预算内覆盖、例外保留和冗余；生成层分别评估正确性、忠实度和引用；端到端层再统计延迟、成本、拒答、权限边界与新鲜度。

在多跳问题中，“命中任意一份相关文档”过于宽松。若每个问题需要若干事实组，每组允许若干替代来源，则证据充分要求每个必要事实组至少有一份有效支持。文档召回后，还需检查选中片段是否包含答案，以及上下文截取后是否保留了必要条件和例外。

定义事件 $R$ 为必要证据被召回，$E$ 为必要证据在组装后仍充分，$G$ 为在充分证据下正确生成。链式法则给出
\begin{equation}
 \Pr(R\cap E\cap G)=\Pr(R)\Pr(E\mid R)\Pr(G\mid R,E).
\end{equation}
该分解不要求三事件独立。若三个条件概率经实测分别为 $0.90$、$0.80$ 与 $0.85$，则受证据严格支持的端到端成功概率为 $0.90\times0.80\times0.85=0.612$。三个条件事件定义于同一请求总体，乘积描述受证据支持的成功路径。缺少证据而凭模型记忆答对的路径需另行统计。

\subsection{引用、忠实与拒答指标}
设可核实声明数为 $M$，受到证据支持的声明数为 $M_s$，可定义声明级忠实度为 $\frac{M_s}{M}$。声明划分规则必须固定，否则模型通过合并或拆分句子就能改变指标。引用精确率按被评估的声明—来源关系中提供支持的比例计算；引用覆盖率衡量需要引用的声明中有多少具有充分引用。两者的分母不同，互换会得到不同含义的数。若支持依赖多份证据的联合，应使用预先声明的集合级判定，要求每一份单独支持全部结论会把联合支持误判为失败。

没有可核实声明的回答，其忠实度分母为零，应单列为不适用；自动记满分会让空回答获得最好成绩。对于不可回答问题，应统计正确拒答率；对于可回答问题，应统计过度拒答率。系统故障和无权限样本应单列，否则接口失败也会被计成拒答成功。模型评审器可扩大检查覆盖面，但需使用人工标注样本估计其误判，并保存版本和判据。

表\ref{tab:rag-evaluation-framework}列出 RAG 常用评测指标及其输入，并给出低分时可检查的环节。组件指标有助于缩小排查范围，具体原因还需通过固定其他组件的对照评估确认。

\begin{table}[htbp]
\centering
\caption{RAG 评测指标与改进方法}
\label{tab:rag-evaluation-framework}
\begin{tabularx}{\linewidth}{@{}lXXXX@{}}
\toprule
评测指标 & 依赖输入与输入对象 & 核心计算逻辑与度量目标 & 指标偏低时的故障归因 & 典型改进策略 \\
\midrule
忠实度 (Faithfulness) & 生成答案、检索上下文 & 答案声明中可被上下文严密支撑的比例 & 生成器内在幻觉、依赖预训练记忆、虚假推理 & 强化证据引用约束、少样本忠实提示、校准拒答阈值 \\
答案相关度 (Answer Relevance) & 用户查询、生成答案 & 答案是否直接切题（常通过逆向生成问题相似度评估） & 答非所问、冗长套话、上下文被无关信息带偏 & 优化提示词指令、强化输出长度与焦点约束 \\
上下文精度 (Context Precision) & 用户查询、检索上下文、标准答案 & 相关文档片段在检索列表中是否优先排在前列 & 重排器精度不足、首部充斥噪声块挤占有效注意力 & 引入交叉编码重排器、优化分块粒度与语义连续性 \\
上下文召回率 (Context Recall) & 用户查询、检索上下文、标准答案 & 标准事实点被检索上下文完整覆盖的比例 & 召回通路漏检、多跳推理依赖缺失、过滤过于激进 & 引入多路混合检索（词法+向量）、查询分解与重写 \\
\bottomrule
\end{tabularx}
\end{table}

\subsection{反事实切片定位故障}
\term{反事实评估}{Counterfactual Evaluation}{}在保持问题与主要配置固定时，改变某个证据条件，观察回答如何变化。它有助于定位依赖关系；相应的因果解释以干预条件与随机性得到控制为前提。

可使用以下成对切片：将检索证据替换为人工确认的充分证据，观察是否仍答错；移除唯一支持某声明的片段，观察是否收缩或拒答；加入条件相同但版本过期的冲突片段，检查时间判断；撤销来源权限，检查检索、上下文和缓存是否共同失效；添加重复、乱序和长干扰片段，检查预算选择的稳定性。还可以固定证据只更换生成器，或固定生成器只更换重排器，使模块差异可解释。

人工充分证据条件用于隔离生成问题，在线评估则同时包含实际检索过程。对于有随机性的生成，应采用相同抽样策略并重复测量；把单个不同的答案一律归因于被改变的模块会高估干预效果。各切片需包含专业术语、跨文档关联、时间限定、低频实体、表格与不可回答问题，防止总体平均值掩盖系统边界。

\subsection{服务效果评估}
若各阶段串行执行，端到端时间近似为排队、改写、检索、重排、组装和生成耗时之和；分支阶段则需按实际依赖图计算关键路径。除了平均值，应报告尾部延迟及超时比例，并说明检索候选数、重排长度、证据词元量与输出长度。增加候选数可能提高召回，同时增加重排成本；增加上下文可能补充证据，也可能挤占输出预算。质量与成本曲线应来自同一配置集合，拼接不同系统的最好结果会画出一条任何单一配置都达不到的曲线。

运行清单使失败可重放，评估集使变化可比较。二者共同回答三个问题：失败发生在哪一层；修复是否改善目标切片；改善是否以更严重的拒答、泄露或超时为代价。只有答案文本而缺少证据版本和模块状态，这类分析就无法完成。



检索、证据选择与生成各阶段的状态记录，可用于重放失败请求并分析召回、支持关系与答案质量。



% BEGIN GENERATED INTERVIEWS

% END GENERATED INTERVIEWS
\chapterreferences
\myendchapterdocument
